\section{Challenges}
\label{sec:challenges}

Extending eBPF programmability to GPUs introduces fundamental challenges arising from (1) monolithic GPU drivers lacking designed-for-programmability interfaces, (2) fundamental mismatches between scalar eBPF semantics and GPU's SIMT execution model, and (3) complexities in coordinating consistent state across asynchronous host-device boundaries.

\subsection{C1: Safe and Efficient Driver Extensibility}

GPU drivers, such as NVIDIA's open GPU kernel modules, are large, monolithic codebases tightly integrating low-level hardware mechanisms and resource-management policies. Unlike Linux kernel subsystems explicitly designed with stable eBPF extension points (e.g., \texttt{struct\_ops}), GPU drivers today do not provide interfaces intended for safe policy programmability. \textbf{Policy logic embedded in driver mechanisms} complicates the introduction of attach points that remain stable across driver versions and GPU generations. For instance, in the UVM kernel module, the page-fault handler function directly embeds LRU-based eviction and tree-based prefetch policies, making it difficult to interpose on memory placement without modifying the driver code path. Moreover, since policy logic executes directly in privileged driver contexts with tight latency constraints and direct GPU hardware control, there are \textbf{strict driver safety requirements}: logic errors (e.g., incorrect lock usage or inconsistent command submission) can cause device-wide stalls or node-level failures, mandating bounded execution and memory safety guarantees.

\subsection{C2: SIMT-Aware Device-Side Execution}

Standard eBPF is designed for scalar, single-threaded execution on CPUs, but GPU kernels operate under the fundamentally different Single Instruction Multiple Threads (SIMT) execution model, creating significant programmability mismatches. Naively mapping scalar eBPF handlers onto GPU threads introduces severe \textbf{warp divergence overhead}: executing scalar policy logic independently on all 32 threads of a warp forces the warp to serialize branches that depend on lane-local state, effectively turning each divergent branch into up to 32 serialized paths in the worst case. Furthermore, GPU architectures feature \textbf{complex memory hierarchies} (per-thread registers, shared memory, L2 caches, global memory) and weak memory consistency semantics. For example, concurrent updates to shared maps from thousands of threads with relaxed atomics cannot be verified by the CPU-side eBPF verifier, which treats programs as single-threaded. Lastly, GPU kernels execute within flat address spaces without isolation or runtime recovery mechanisms; policy errors like infinite loops or invalid memory accesses cause unrecoverable hardware stalls, making \textbf{bounded execution and memory safety guarantees} essential to enforce statically at compile time.

\subsection{C3: Shared Cross-Layer State and Coordination}

Many GPU policies require coordinated actions spanning CPU host drivers and GPU device execution; for example, GPU-side instrumentation capturing warp-level memory accesses informs host-side memory placement and scheduling decisions. Ensuring efficient and consistent cross-layer coordination poses unique technical difficulties. First, there is a significant \textbf{cross-boundary memory hierarchy mismatch}: placing shared maps naively in host memory introduces prohibitive latency (tens to hundreds of microseconds per access from GPU warps via PCIe), drastically impacting GPU utilization, whereas placing all shared state within GPU global memory risks limited capacity, bandwidth bottlenecks, and cache thrashing. Second, host driver operations and GPU execution proceed asynchronously at vastly different granularities, complicating \textbf{state consistency}. Without careful coordination, host-side policy decisions (e.g., page eviction or kernel scheduling) risk relying on stale GPU-side signals (e.g., warp-level memory access statistics), causing thrashing migrations and oscillatory scheduling. In \S\ref{sec:verification} we describe how \sys addresses cross-layer state management.

\section{\sys Design}

\subsection{Design Principles}

\sys follows three fundamental design principles:

\paragraph{Principle 1: Transparent and Dynamic Policy Programmability.}
\sys explicitly separates stable, privileged GPU driver mechanisms from dynamically programmable policies. Instead of embedding policies in opaque vendor stacks or application-specific runtimes, \sys provides clearly-defined policy interfaces that allow dynamic updates without requiring driver modifications or invasive application changes.

\paragraph{Principle 2: Unified Policy State and Control Plane.}
To eliminate the complexity and fragmentation of heterogeneous GPU stacks, \sys offers a unified abstraction for policy state, programming interfaces, and lifecycle management across CPU hosts and GPU devices. This ensures consistent cross-layer coordination, simplifies policy expression, and enables coherent decisions in asynchronous execution environments.

\paragraph{Principle 3: Verifiable, SIMT-aware Safe and Efficient Execution.}
Recognizing the fundamental mismatch between traditional scalar eBPF semantics and GPU's SIMT parallelism, \sys introduces a static verification model explicitly designed for SIMT execution. Through a type discipline that distinguishes warp-uniform from lane-varying values, the verifier statically enforces bounded resource usage, latency constraints, and memory safety, eliminating the need for runtime isolation or overhead.

\subsection{High-Level Architecture}
\label{sec:arch}

\begin{figure}
\centering
\fbox{\parbox{0.8\columnwidth}{\centering [Figure: \sys High-Level Design]}}
\vspace{-.5ex}
\caption{\sys high-level design: unified eBPF runtime spanning CPU and GPU.}
\label{fig:system-design}
\end{figure}

Figure~\ref{fig:system-design} shows the high-level architecture of \sys, composed of a user-space control plane and a unified kernel-device execution path. \sys policies form distributed applications consisting of host- and device-side eBPF programs, hierarchical cross-layer maps, and a user-space control-plane for lifecycle management.

\paragraph{User-space Control Plane.}
Policy authors use standard eBPF tooling (e.g., \texttt{clang/libbpf}, \texttt{bpftool}) to build policy bytecode. The control-plane links against a loader library that installs policies and configures hierarchical maps. It enables runtime policy redeployment and reconfiguration (e.g., eviction thresholds, sampling rates, priorities) without application or kernel restarts, and exports runtime metrics (e.g., region hotness, tenant latencies) to external systems.

\paragraph{Kernel and Device Runtime.}
\sys provides a unified runtime across kernel and GPU contexts, managing verification, compilation, and map placement. It employs a SIMT-aware verifier (\S\ref{sec:verification}) that extends standard eBPF checks (warp-uniform semantics, bounded execution, GPU synchronization constraints). Verified programs are JIT-compiled into native host code or GPU-compatible instructions (e.g., PTX). The runtime transparently manages hierarchical cross-layer maps, automatically optimizing their placement and synchronization across host memory and GPU memory tiers. Policy handlers access maps exclusively via standard eBPF helpers, eliminating explicit cross-layer synchronization or address translation.

Policy execution occurs at structured hook points defined in GPU drivers and kernels. On the host, \sysdriver exposes \texttt{struct\_ops}-style hooks (e.g., region admission/eviction, kernel scheduling). Host handlers return concise policy decisions enforced via safe kernel helpers. On the GPU device, \sys injects lightweight trampolines at kernel entry/exit, warp synchronization barriers, and critical memory operations, invoking device-side handlers at warp granularity. Device handlers make granular control decisions (e.g., dynamic prefetching, block-level scheduling) and update cross-layer maps with detailed performance signals (e.g., region access patterns, warp occupancy). They never directly manipulate hardware registers or page tables.

Together, the control-plane, unified runtime, and structured hook execution form an adaptive closed loop: device-side logic actively manages GPU resources and instrumentation; host-side handlers enforce global policies using shared instrumentation state; and the control-plane dynamically monitors and updates policies in real-time.

\subsection{Policy Interfaces}
\label{sec:design-interfaces}

\sys introduces two primary eBPF program types to categorize policy logic:
a driver-level type for GPU policy hooks (encompassing both \texttt{BPF\_PROG\_TYPE\_GPU\_MEM} for memory events and \texttt{BPF\_PROG\_TYPE\_GPU\_SCHED} for scheduling) and a device-level type (\texttt{BPF\_PROG\_TYPE\_GPU\_DEV}) for GPU-kernel instrumentation. These types coexisting with existing kernel eBPF program types (e.g., cgroup, TC) without conflict.

\subsubsection{Memory Interface}
\label{sec:design-mem}

\sys treats memory placement as a programmable region
cache: it exposes a small set of region-level events for \emph{admission},
\emph{access}, and \emph{removal}, together with a prefetch channel, and lets host
and device cooperate on a unified policy.

On the host driver, these events are realized as a small \texttt{struct\_ops}-style policy
table with typed contexts:

\begin{minted}[highlightlines=false,fontsize=\small]{c}
// Host-side policy callbacks for GPU memory management
struct gdrv_mem_ops {
  // Region created or eligible for device placement
  int (*region_add)(gdrv_mem_add_ctx_t *ctx);
  // Memory subsystem observes region use; decide cache hit
  int (*region_access)(gdrv_mem_access_ctx_t *ctx);
  // Region about to leave device-resident working set
  int (*region_remove)(gdrv_mem_remove_ctx_t *ctx);
  // Safe points (kernel launches, page faults) for prefetch
  int (*prefetch)(struct gdrv_mem_prefetch_ctx_t *ctx);
};
// Host-side kfunc: operate on region (VA unit) across GPU/CPU
bool gdrv_mem_list_insert_head(gdrv_mem_list_t* head, 
                               gmem_region_t* r);
bool gdrv_mem_list_insert_tail(gdrv_mem_list_t* tail, 
                               gmem_region_t* r);

// Device-side memory callback interface
struct gdev_mem_ops {
    void (*memtrace)(...);
    void (*memsync)(...);
};
// Device-side helpers
// Request prefetch, triggers prefetch callback in kernel
int gdev_mem_prefetch_hint(struct gmem_region* r);
// Pin memory on CPU or GPU
int gdev_mem_pin_hint(struct gmem_region* r, int device_id);
\end{minted}

\paragraph{Host-side region events and removal.}
On the host, \sys models memory as a set of logical regions and exposes attach points for three
canonical cache events. \sys maintains all device-resident regions in a kernel-owned doubly linked list. The memory subsystem always inserts newly admitted regions at the tail and evicts from the head when under pressure, ensuring eventual eviction of all regions even if policies are buggy. eBPF handlers may only reorder regions via \texttt{gdrv\_mem\_list\_*} helpers, and cannot remove the head node; the kernel retains authority to evict head regions unconditionally.

Together, these events realize a
general region-cache interface: policies implemented as eBPF handlers maintain their
own state in shared maps (for example, LRU lists, LFU counters, or tenant budgets),
and respond to Add/Access/Remove and prefetch callbacks with compact decisions
such as \emph{admit}, \emph{bypass}, \emph{pin}, \emph{evict}, or ordered lists
of regions to migrate. The underlying memory subsystem enforces these decisions
and maintains correct mappings, but delegates replacement and prefetch policy to
\sys.

\paragraph{device-side memory events.}

On the GPU, \sys exposes a single device-side memory event interface that kernels
use both to report observed accesses and to express future-use hints. Device-side
handlers receive a compact context struct (not shownp') with fields such as a logical
region identifier, kernel, warp and SM identifiers, an event type (\emph{ACCESS}
or \emph{HINT}), and optional hint flags.

At selected memory operation sites, \sys injects lightweight hooks that operate at
warp granularity and invoke a device-side \sys handler with an \emph{ACCESS} event:
the handler updates shared maps with per-region statistics such as access counts,
approximate working-set membership, and simple locality signatures.\

The same device-side interface also supports \emph{HINT} events, which kernels
invoke explicitly when they have domain knowledge about future region usage. For
example, a mixture-of-experts kernel can compute the identifiers of experts it
will access for the next batch and call a \sys helper with HINT events that mark
the corresponding regions as "will-use soon" with an optional priority or deadline;
a KV-cache kernel can emit HINT events that mark segments as cold and suitable for
demotion. Both ACCESS and HINT events are recorded in shared maps or small control
queues that host-side \sys handlers consult at Access, Remove, and prefetch hooks.
Device-side code never manipulates page tables or performs migration directly; it
implements the sensing and hinting side of a closed-loop policy, while host-side
\sys handlers implement the placement, prefetch, and eviction decisions over the
GPU memory hierarchy.


\subsubsection{Scheduling Interface}
\label{sec:design-sched}

GPU schedulers today typically implement fixed, coarse-grained policies such as FIFO
queues with limited priority control, which are insufficient in multi-tenant settings
with mixed latency-sensitive and batch workloads or tight power constraints. In such
environments, operators need policies that can shape both \emph{which kernels reach
the GPU} at the granularity of kernel, and \emph{how work is consumed inside those kernels} at the granularity of thread blocks or tiles. \sys exposes scheduling interfaces on both host and device
so that policies can implement fine-grained priority, deadline-aware, and power-aware
scheduling while respecting the GPU SIMT execution model.

On the host, these lifecycle events are exported as a small struct-ops table of
scheduling callbacks:

\begin{minted}[highlightlines=false,fontsize=\small]{c}
// Host-side policy callbacks for GPU kernel scheduling
struct gdrv_sched_ops {
  // Kernel enqueued; return 0 to accept, <0 to reject/defer
  int  (*submit)(gsched_submit_ctx_t *ctx);
  // Dispatch ready kernel; return kernel ID or -1 to defer
  int  (*admit)(gsched_admit_ctx_t *ctx);
  // Kernel completed notification; must not fail
  void (*complete)(gsched_complete_ctx_t *ctx);
};
// Host-side kfunc
bool gdrv_sched_preempt(gdrv_preempt_ctx_t *ctx);
// Device-side block scheduling callbacks
struct gdev_sched_ops {
  // Worker block starts new work unit
  void (*entry)(gdev_block_ctx_t *ctx);
  // Worker block finishes work unit
  void (*exit)(gdev_block_ctx_t *ctx);
  // return whether to steal work (Thread Block scheduler)
  bool (*should_try_steal)(gdev_block_ctx_t *ctx);
};
\end{minted}

\paragraph{Host-side scheduling hooks.}
On the host, \sys extends the GPU kernel lifecycle with attach points that allow
policies to control admission, ordering, and preemption of kernels:

Attached eBPF schedulers read these contexts and shared maps (for example, per-tenant
latency and throughput statistics, outstanding work across queues) and return compact
decisions such as which kernel to run next, whether to throttle a tenant, or whether
to allow preemption. The host-side interface thus controls which kernels occupy GPU
contexts and how they are multiplexed in multi-tenant scenarios.
\paragraph{Device-side thread-block scheduler and signals.}

Inside the GPU, \sys exposes a complementary interface that operates at the level of
thread blocks or tiles, aligning with the SIMT execution model. Each kernel that uses \sys's internal scheduler is modeled as a pool of \emph{work units} (e.g., tiles or logical block indices) consumed by \emph{worker blocks}. Device-side policies observe and steer this internal scheduler via a \texttt{gdev\_block\_ctx} that includes the kernel identifier, tenant identifier, and a logical work-unit or block identifier.

Together, the host-side and device-side interfaces allow \sys policies to coordinate decisions across the full scheduling pipeline: from kernel admission and preemption on the host to fine-grained block-level work distribution inside GPU kernels.


\subsection{Runtime Verification and Optimizations}
\label{sec:verification}

\sys policies execute directly within privileged GPU driver and kernel contexts, making strong safety guarantees essential. We adopt a standard eBPF threat model: we trust cluster operators loading policies, but consider policy code potentially buggy or misconfigured. The trusted computing base (TCB) includes the kernel eBPF verifier, \sys-specific kernel module, GPU runtime backend, and GPU driver/firmware. Policy authors only interact with eBPF abstractions and do not write native device code. For host-side verification, we reuse the existing Linux kernel eBPF verifier unchanged, relying on standard type-checking, bounded-loop, and memory-safety constraints. \sys-specific kfuncs and hooks are declared via BPF Type Format (BTF) metadata.

\subsubsection{SIMT-aware Verification}

A fundamental mismatch exists between traditional scalar eBPF semantics and GPU's SIMT parallel execution model. Directly executing scalar eBPF handlers on GPUs would severely degrade performance due to increased warp divergence and unbounded resource consumption. To address this, \sys introduces a novel SIMT-aware static verification model that rigorously enforces warp-level uniformity constraints and guarantees bounded resource usage.

The core insight is explicitly distinguishing \emph{warp-uniform} values (constant across all threads within a warp) from \emph{lane-varying} values (different per thread). We define a new type discipline integrated into the standard kernel eBPF verifier. The verifier enforces that critical control-flow decisions (branch conditions, loop bounds) strictly depend on warp-uniform variables. This constraint ensures that warp threads follow identical control paths, enabling safe and efficient warp-level aggregated execution.

Our verification model statically ensures safety-critical properties: bounded execution, no infinite loops or stalls, no unauthorized memory accesses, and freedom from prohibited synchronization operations. This explicit, GPU-specific extension of the eBPF verification model allows safe injection of programmable policies directly into privileged GPU driver and kernel contexts.

\subsubsection{Warp-level Execution Optimizations}

While SIMT-aware verification ensures correctness and safety, direct scalar execution of eBPF handlers across GPU threads would still incur prohibitive divergence and memory overhead. To eliminate this inefficiency, \sys introduces a novel warp-level aggregated execution model optimized specifically for GPU hardware characteristics.

The key idea is executing scalar policy logic exactly once per warp, using a designated ``warp leader'' thread. Warp-level primitives (e.g., shuffle, ballot) efficiently gather lane-varying inputs (such as per-thread addresses or activity masks) from all threads. The warp leader computes policy decisions based on these aggregated inputs, and broadcasts compact decisions back to the warp. Under warp-uniform constraints enforced by the verifier, this aggregated warp-level execution preserves the original scalar eBPF semantics while drastically minimizing GPU overhead.

Additionally, recognizing that GPU policy code frequently accesses local, highly-thread-local map state, we perform aggressive specialization and inlining optimizations. These compiler-level optimizations eliminate unnecessary branching and helper call overhead, further reducing instruction count and warp divergence.

\subsubsection{Hierarchical Cross-layer Maps for State Coordination}

Effective GPU policy decisions require efficient, coherent sharing of policy state across asynchronous CPU and GPU execution contexts. Traditional eBPF maps are CPU-centric and lack explicit cross-layer coordination semantics, causing unacceptable performance penalties if naively extended to GPUs. To address this, \sys introduces a novel hierarchical cross-layer map abstraction explicitly designed to transparently bridge CPU and GPU memory hierarchies.

The central design insight is to provide developers a unified logical abstraction of key-value maps that transparently spans multiple memory tiers: CPU host DRAM, GPU global memory, and GPU on-chip memories. Policy developers never manually handle explicit DMA transfers, memory placements, synchronization, or address translations. Internally, the runtime dynamically manages map placement and synchronization based on latency, capacity, and bandwidth constraints, automatically balancing locality and coherence needs.

This abstraction allows developers to express complex cross-layer logic (e.g., GPU-side region access tracking informing host-side eviction policies) using simple eBPF APIs without explicit synchronization overhead.
